\frfilename{mt275.tex}
\versiondate{3.12.12}

\def\chaptername{Teori probabilitas}
\def\sectionname{Martingal}
\copyrightdate{2001}

\newsection{275}

Sejauh ini bab ini didominasi oleh barisan variabel acak yang saling bebas.
Sekarang saya beralih pada satu lagi konsep luar biasa yang lahir dari intuisi
probabilistik. Di sini kita mempelajari sistem yang berkembang, tempat kita
memperoleh semakin banyak informasi seiring berjalannya waktu. Saya memberikan
teorema-teorema dasar tentang konvergensi titik demi titik martingal
(275F-275H, %275F 275G 275H
275K) dan uraian yang sangat singkat tentang `waktu henti'
(275L-275P). %275L 275M 275N 275O 275P

\leader{275A}{Definisi} Misalkan $(\Omega,\Sigma,\mu)$ adalah ruang
probabilitas dengan pelengkapan $(\Omega,\hat\Sigma,\hat\mu)$, dan
$\sequencen{\Sigma_n}$ barisan subaljabar-sigma dari $\hat\Sigma$ yang tak
menurun. \cmmnt{(Barisan $\sequencen{\Sigma_n}$ demikian disebut
{\bf filtrasi}.)}
Suatu {\bf martingal yang teradaptasi terhadap
$\sequencen{\Sigma_n}$} adalah barisan $\sequencen{X_n}$ dari variabel acak
bernilai real yang terintegralkan pada $\Omega$ sedemikian sehingga
(i) $\dom X_n\in\Sigma_n$ dan $X_n$ terukur terhadap $\Sigma_n$ untuk setiap
$n\in\Bbb N$ (ii) kapan pun $m\le n\in\Bbb N$ dan $E\in\Sigma_m$, maka
$\int_EX_{n}=\int_EX_m$.

\cmmnt{Perhatikan bahwa untuk (ii) cukup jika
$\int_EX_{n+1}=\int_EX_n$ kapan pun $n\in\Bbb N$ dan $E\in\Sigma_n$.}

\leader{275B}{Contoh }\cmmnt{Kita telah melihat banyak konteks tempat
barisan semacam ini muncul secara alami; berikut beberapa di antaranya.

\medskip

}{\bf (a)} Misalkan $(\Omega,\Sigma,\mu)$ adalah ruang probabilitas dan
$\sequencen{\Sigma_n}$ barisan subaljabar-sigma dari $\Sigma$ yang tak
menurun. Misalkan $X$ adalah sebarang variabel acak bernilai real pada
$\Omega$ dengan ekspektasi berhingga, dan untuk setiap $n\in\Bbb N$ misalkan
$X_n$ adalah suatu ekspektasi bersyarat dari $X$ pada $\Sigma_n$\cmmnt{,
seperti dalam \S233}. Dengan tunduk pada syarat bahwa
$\dom X_n\in\Sigma_n$ dan $X_n$ benar-benar terukur terhadap $\Sigma_n$
untuk setiap $n$\cmmnt{ (suatu hal yang murni teknis -- lihat 232He)},
$\sequencen{X_n}$ merupakan martingal yang teradaptasi terhadap
$\sequencen{\Sigma_n}$\cmmnt{, karena
$\int_EX_{n+1}=\int_EX=\int_EX_n$ kapan pun $E\in\Sigma_n$}.

\spheader 275Bb Misalkan $(\Omega,\Sigma,\mu)$ adalah ruang probabilitas
dan $\sequencen{X_n}$ barisan variabel acak yang saling bebas dan semuanya
berekspektasi nol. Untuk setiap $n\in\Bbb N$ misalkan
$\tilde\Sigma_n$ adalah aljabar-sigma yang dibangkitkan oleh
$\bigcup_{i\le n}\Sigma_{X_i}$, dengan $\Sigma_{X_i}$ menyatakan
aljabar-sigma yang didefinisikan oleh $X_i$\cmmnt{ (272C)}, dan tetapkan
$S_n=X_0+\ldots+X_n$. Maka $\sequencen{S_n}$ adalah martingal yang
teradaptasi terhadap $\tilde\Sigma_n$. \cmmnt{(Gunakan 272K untuk melihat
bahwa $\Sigma_{X_{n+1}}$ bebas dari $\tilde\Sigma_n$, sehingga
$\int_EX_{n+1}=\int X_{n+1}\times\chi E=0$ untuk setiap
$E\in\tilde\Sigma_n$, berdasarkan 272R.)}

\spheader 275Bc Misalkan $(\Omega,\Sigma,\mu)$ adalah ruang probabilitas
dan $\sequencen{X_n}$ barisan variabel acak yang saling bebas dan semuanya
berekspektasi $1$. Untuk setiap $n\in\Bbb N$ misalkan
$\tilde\Sigma_n$ adalah aljabar-sigma yang dibangkitkan oleh
$\bigcup_{i\le n}\Sigma_{X_i}$, dengan $\Sigma_{X_i}$ menyatakan
aljabar-sigma yang didefinisikan oleh $X_i$, dan tetapkan
$W_n=X_0\times\ldots\times X_n$. Maka $\sequencen{W_n}$ adalah martingal
yang teradaptasi terhadap $\sequencen{\tilde\Sigma_n}$.
%\sequencen{\tilde\Sigma_n} is the {\bf natural filtration} for
%\sequencen{X_n}

\leader{275C}{Catatan }\cmmnt{{\bf (a)} Tampaknya sesuai dengan konsep
variabel acak $X$ yang `teradaptasi' terhadap aljabar-sigma $\Sigma$ untuk
mensyaratkan bahwa $\dom X\in\Sigma$ dan $X$ terukur terhadap $\Sigma$,
meskipun ini dapat berarti bahwa variabel acak lain, yang sama hampir di
mana-mana dengan $X$, mungkin tidak `teradaptasi' terhadap $\Sigma$.

\spheader 275Cb Masalah teknis semacam ini tentu lenyap jika semua himpunan
yang terabaikan terhadap $\mu$ di dalam $\Omega$ termasuk dalam $\Sigma_0$.
Namun contoh seperti 275Bb membuatnya tampak tidak wajar untuk memaksakan
penyederhanaan semacam ini sebagai aturan umum.

\header{275Cc}}%end of comment
{\bf (c)} Konsep `martingal' dapat dengan mudah diperluas ke himpunan indeks
selain $\Bbb N$; bahkan, jika $I$ adalah sebarang himpunan terurut parsial,
kita dapat mengatakan bahwa $\langle X_i\rangle_{i\in I}$ adalah martingal
pada $(\Omega,\Sigma,\mu)$ yang teradaptasi terhadap
$\langle\Sigma_i\rangle_{i\in I}$ jika (i) setiap $\Sigma_i$ adalah
subaljabar-sigma dari $\hat\Sigma$ (ii) setiap $X_i$ adalah variabel acak
bernilai real yang terintegralkan dan terukur terhadap $\Sigma_i$, dengan
$\dom X_i\in\Sigma_i$ (iii) kapan pun $i\le j$ dalam $I$, maka
$\Sigma_i\subseteq\Sigma_j$ dan $\int_EX_i=\int_EX_j$ untuk setiap
$E\in\Sigma_i$. \cmmnt{Kasus utama setelah $I=\Bbb N$ adalah
$I=\coint{0,\infty}$; $I=\Bbb Z$ juga menarik, dan saya rasa wajar untuk
mengatakan bahwa gagasan-gagasan terpenting sudah dapat diungkapkan dalam
teorema tentang martingal yang diindeks oleh himpunan berhingga $I$. Namun
dalam jilid ini pada umumnya saya akan mengambil martingal yang diindeks oleh
$\Bbb N$.
}%end of comment

\spheader 275Cd Jika hanya diberikan barisan $\sequencen{X_n}$ dari
variabel acak bernilai real yang terintegralkan pada ruang probabilitas
$(\Omega,\Sigma,\mu)$, kita cukup mengatakan bahwa $\sequencen{X_n}$ adalah
suatu {\bf martingal} pada $(\Omega,\Sigma,\mu)$ jika terdapat suatu barisan
subaljabar-sigma $\sequencen{\Sigma_n}$ dari $\hat\Sigma$ yang tak menurun
\cmmnt{ (pelengkapan $\Sigma$)} sedemikian sehingga $\sequencen{X_n}$ adalah
martingal yang teradaptasi terhadap $\sequencen{\Sigma_n}$. Jika kita
menuliskan $\tilde\Sigma_n$ untuk aljabar-sigma yang dibangkitkan oleh
$\bigcup_{i\le n}\Sigma_{X_i}$, dengan $\Sigma_{X_i}$ adalah aljabar-sigma
yang didefinisikan oleh $X_i$\cmmnt{, seperti dalam 275Bb}, maka\cmmnt{
mudah dilihat bahwa} $\sequencen{X_n}$ adalah martingal jika dan hanya jika
ia adalah martingal yang teradaptasi terhadap
$\sequencen{\tilde\Sigma_n}$.

\spheader 275Ce Melanjutkan dari (d),\cmmnt{ mudah pula dilihat bahwa} jika
$\sequencen{X_n}$ adalah martingal pada $(\Omega,\Sigma,\mu)$, dan
$X'_n\eae X_n$ untuk setiap $n$, maka $\sequencen{X'_n}$ adalah martingal
pada $(\Omega,\Sigma,\mu)$. \cmmnt{(Intinya ialah bahwa jika
$\sequencen{X_n}$ teradaptasi terhadap $\sequencen{\Sigma_n}$, maka baik
$\sequencen{X_n}$ maupun $\sequencen{X'_n}$ teradaptasi terhadap
$\sequencen{\hat\Sigma_n}$, dengan

\Centerline{$\hat\Sigma_n
=\{E\symmdiff F:E\in\Sigma_n,\,F\text{ terabaikan}\}$.)}

\noindent}Akibatnya kita mempunyai konsep `martingal' sebagai barisan dalam
$L^1(\mu)$, yaitu suatu barisan $\sequencen{X_n^{\ssbullet}}$ dalam
$L^1(\mu)$ adalah martingal jika dan hanya jika $\sequencen{X_n}$ adalah
martingal.

\cmmnt{Meskipun demikian, saya berpendapat bahwa konsep `martingal yang
teradaptasi terhadap suatu barisan aljabar-sigma' adalah konsep utama, karena
dalam semua penerapan pokok aljabar-sigma tersebut mencerminkan suatu aspek
esensial dari masalah yang mungkin tidak sepenuhnya tercakup oleh variabel acak
saja.}

\cmmnt{
\spheader 275Cf Kata `martingale' mula-mula (dalam bahasa Inggris; sejarahnya
dalam bahasa Prancis lebih rumit) merujuk pada tali kekang yang digunakan
untuk mencegah kuda mendongakkan kepalanya. Kemudian kata ini digunakan untuk
menamai sistem perjudian tempat penjudi menggandakan taruhannya setiap kali
kalah, dan (dalam bahasa Prancis) sebagai istilah umum bagi sistem perjudian.
Sistem-sistem ini dapat dipandang sebagai suatu kelas `martingal waktu henti',
seperti dijelaskan dalam 275L-275P %275L 275M 275N 275O 275P
di bawah.
}%end of comment

\leader{275D}{}\cmmnt{ Sebagian besar teori martingal terdiri atas berbagai
jenis ketaksamaan. Saya memberikan dua yang paling penting, keduanya berasal
dari J.L.Doob. (Lihat pula 276Xa-276Xb.)

\medskip

\noindent}{\bf Lemma} Misalkan $(\Omega,\Sigma,\mu)$ adalah ruang
probabilitas, dan $\sequencen{X_n}$ suatu martingal pada $\Omega$. Tetapkan
$n\in\Bbb N$ dan $X^*=\max(X_0,\ldots,X_n)$. Maka untuk setiap
$\epsilon>0$,

\Centerline{$\Pr(X^*\ge\epsilon)
\le\Bover1{\epsilon}\Expn(X_n^+)$,}

\noindent dengan $X_n^+=\max({0},X_n)$.

\proof{ Tuliskan $\hat\mu$ untuk pelengkapan $\mu$, dan $\hat\Sigma$ untuk
domainnya. Misalkan $\sequencen{\Sigma_n}$ adalah barisan subaljabar-sigma
dari $\hat\Sigma$ yang tak menurun dan terhadapnya $\sequencen{X_n}$
teradaptasi. Untuk setiap $i\le n$ tetapkan

\Centerline{$E_i=\{\omega:\omega\in\dom
X_i,\,X_i(\omega)\ge\epsilon\}$,}

\Centerline{$F_i=E_i\setminus\bigcup_{j<i}E_j$.}

\noindent Maka $F_0,\ldots,F_n$ saling lepas dan
$F=\bigcup_{i\le n}F_i=\bigcup_{i\le n}E_i$; selain itu, dengan $H$
menyatakan himpunan koterabaikan $\bigcap_{i\le n}\dom X_i$,

\Centerline{$\{\omega:X^*(\omega)\ge\epsilon\}=F\cap H$,}

\noindent sehingga

\Centerline{$\Pr(X^*\ge\epsilon)=
\hat\mu\{\omega:X^*(\omega)\ge\epsilon\}
=\hat\mu F=\sum_{i=0}^n\hat\mu F_i$.}

\noindent Di sisi lain, $E_i$ dan $F_i$ termasuk dalam $\Sigma_i$ untuk
setiap $i\le n$, sehingga

\Centerline{$\int_{F_i}X_n=\int_{F_i}X_i\ge \epsilon\hat\mu F_i$}

\noindent untuk setiap $i$, dan

\Centerline{$\epsilon\hat\mu F=\epsilon\sum_{i=0}^n\hat\mu F_i
\le\sum_{i=0}^n\int_{F_i}X_n
=\int_FX_n
\le\int_FX_n^+
\le\Expn(X_n^+)$,}

\noindent sebagaimana diperlukan.
}%end of proof of 275D

\cmmnt{\medskip

\noindent{\bf Catatan} Perhatikan bahwa sebenarnya kita mempunyai
$\epsilon\hat\mu F\le\int_FX_n$, dengan
$F=\{\omega:X^*(\omega)\ge\epsilon\}$; hal ini sangat penting dalam banyak
penerapan.
}%end of comment

\leader{275E}{Pelintasan naik}\cmmnt{ Lemma berikut bergantung pada gagasan
`pelintasan naik'.} Misalkan $x_0,\ldots,x_n$ adalah sebarang daftar bilangan
real, dan $a<b$ dalam $\Bbb R$. {\bf Banyaknya pelintasan naik dari $a$ ke
$b$} dalam daftar $x_0,\ldots,x_n$ adalah banyaknya pasangan $(j,k)$
sedemikian sehingga $0\le j<k\le n$, $x_j\le a$, $x_k\ge b$, dan
$a<x_i<b$ untuk $j<i<k$. \cmmnt{Perhatikan bahwa ini juga merupakan $m$
terbesar sedemikian sehingga $s_m<\infty$, jika kita menulis

\Centerline{$r_1=\inf\{i:i\le n,\,x_i\le a\}$,}

\Centerline{$s_1=\inf\{i:r_1<i\le n,\,x_i\ge b\}$,}

\Centerline{$r_2=\inf\{i:s_1<i\le n,\,x_i\le a\}$,}

\Centerline{$s_2=\inf\{i:r_2<i\le n,\,x_i\ge b\}$}

\noindent dan seterusnya, dengan mengambil $\inf\emptyset=\infty$.
}%end of comment

\leader{275F}{Lemma} Misalkan $(\Omega,\Sigma,\mu)$ adalah ruang
probabilitas dan $\sequencen{X_n}$ suatu martingal pada $\Omega$. Anggap bahwa
$n\in\Bbb N$ dan $a<b$ dalam $\Bbb R$. Untuk setiap
$\omega\in\bigcap_{i\le n}\dom X_i$, misalkan $U(\omega)$ adalah banyaknya
pelintasan naik dari $a$ ke $b$ dalam daftar
$X_0(\omega),\ldots,X_n(\omega)$. Maka

\Centerline{$\Expn(U)\le\Bover1{b-a}\Expn((X_n-X_0)^+)$,}

\noindent dengan $(X_n-X_0)^+(\omega)=\max(0,X_n(\omega)-X_0(\omega))$
untuk $\omega\in\dom X_n\cap\dom X_0$.

\proof{ Setiap langkah tersendiri dalam pembuktian ini bersifat `elementer',
tetapi strukturnya secara keseluruhan tidak sepele.

\medskip

{\bf (a)} Fakta berikut akan berguna. Misalkan $x_0,\ldots,x_n$ adalah
bilangan real; misalkan $u$ adalah banyaknya pelintasan naik dari $a$ ke $b$
dalam daftar $x_0,\ldots,x_n$. Tetapkan $y_i=\max(x_i,a)$ untuk setiap $i$;
maka $u$ juga merupakan banyaknya pelintasan naik dari $a$ ke $b$ dalam daftar
$y_0,\ldots,y_n$. Untuk setiap $k\le n$, tetapkan $c_k=1$ jika terdapat
$j\le k$ sedemikian sehingga $x_j\le a$ dan $x_i<b$ untuk $j\le i\le k$,
dan $0$ jika tidak. Maka

\Centerline{$(b-a)u\le\sum_{k=0}^{n-1}c_k(y_{k+1}-y_{k})$.}

\noindent\vthsp\Prf\ Saya melakukan induksi pada $m$ untuk menunjukkan bahwa
(dengan mendefinisikan $r_m$, $s_m$ seperti dalam 275E)

\Centerline{$(b-a)m\le\sum_{k=0}^{s_m-1}c_k(y_{k+1}-y_{k})$}

\noindent kapan pun $m\le u$. Untuk $m=0$ (dengan mengambil $s_0=-1$) kita
mempunyai $0=0$. Untuk langkah induksi menuju $m\ge 1$, kita mempunyai
$s_{m-1}<r_{m}<s_m\le n$ (karena saya mengandaikan $m\le u$), dan
$c_k=0$ jika $s_{m-1}\le k<r_{m}$, $c_k=1$ jika $r_{m}\le k<s_m$. Jadi

$$\eqalignno{\sum_{k=0}^{s_{m}-1}c_k(y_{k+1}-y_{k})
&=\sum_{k=0}^{s_{m-1}-1}c_k(y_{k+1}-y_{k})
+\sum_{k=r_m}^{s_m-1}(y_{k+1}-y_k)\cr
&\ge(b-a)(m-1)+y_{s_m}-y_{r_m}\cr
\noalign{\noindent (berdasarkan hipotesis induksi)}
&\ge(b-a)m\cr}$$

\noindent (karena $y_{s_m}\ge b$, $y_{r_m}=a$), dan induksi berlanjut.


Dengan demikian

\Centerline{$\sum_{k=0}^{s_u-1}c_k(y_{k+1}-y_{k})\ge(b-a)u$.}

\noindent Untuk jumlah $\sum_{k=s_u}^{n-1}c_k(y_{k+1}-y_k)$, kita mempunyai
$c_k=0$ untuk $s_u\le k<r_{u+1}$, $c_k=1$ untuk
$r_{u+1}\le k<s_{u+1}$, sedangkan $s_{u+1}>n$, sehingga jika
$n\le r_{u+1}$ kita mempunyai

\Centerline{$\sum_{k=0}^{n-1}c_k(y_{k+1}-y_{k})
=\sum_{k=0}^{s_u-1}c_k(y_{k+1}-y_{k})\ge(b-a)u$,}

\noindent sedangkan jika $n>r_{u+1}$ kita mempunyai

$$\eqalign{\sum_{k=0}^{n-1}c_k(y_{k+1}-y_{k})
&=\sum_{k=0}^{s_u-1}c_k(y_{k+1}-y_{k})
+\sum_{k=r_{u+1}}^{n-1}y_{k+1}-y_k\cr
&\ge(b-a)u+y_n-y_{r_{u+1}}\cr
&\ge(b-a)u\cr}$$


\noindent karena $y_n\ge a=y_{r_{u+1}}$. Jadi dalam kedua kasus kita
memperoleh hasil yang diperlukan.\ \Qed

\medskip

{\bf (b)(i)} Sekarang definisikan

\Centerline{$Y_k(\omega)=\max(a,X_k(\omega))$ untuk $\omega\in\dom X_k$,}

\Centerline{$F_k=\{\omega:\omega\in\bigcap_{i\le k}\dom X_i,\,
\exists\,j\le k,\,X_j(\omega)\le a,\,X_i(\omega)<b$ jika $j\le i\le k\}$}

\noindent untuk setiap $k\in\Bbb N$. Jika $\sequencen{\Sigma_n}$ adalah
barisan aljabar-sigma tak menurun yang terhadapnya $\sequencen{X_n}$
teradaptasi, maka $F_k\in\Sigma_k$ (karena jika $j\le k$, semua himpunan
$\dom X_j$, $\{\omega:X_j(\omega)\le a\}$,
$\{\omega:X_j(\omega)<b\}$ termasuk dalam
$\Sigma_j\subseteq\Sigma_k$).

\medskip

\quad{\bf (ii)} Kita mendapati bahwa $\int_FY_k\le\int_FY_{k+1}$ jika
$F\in\Sigma_k$. \Prf\ Tetapkan $G=\{\omega:X_k(\omega)>a\}\in\Sigma_k$.
Maka

$$\eqalign{\int_FY_k
&=\int_{F\cap G}X_k+a\hat\mu(F\setminus G)\cr
&=\int_{F\cap G}X_{k+1}+a\hat\mu(F\setminus G)\cr
&\le\int_{F\cap G}Y_{k+1}+\int_{F\setminus G}Y_{k+1}
=\int_FY_{k+1}.\,\,\Qed\cr}$$

\medskip

\quad{\bf (iii)}
Akibatnya $\int_FY_{k+1}-Y_k\le\int Y_{k+1}-Y_k$ untuk setiap
$F\in\Sigma_k$.

\Centerline{\Prf\enskip $\int (Y_{k+1}-Y_k)-\int_F(Y_{k+1}-Y_k)
=\int_{\Omega\setminus F}Y_{k+1}-\int_{\Omega\setminus F}Y_k\ge 0$.
\Qed}

\medskip

{\bf (c)} Misalkan $H$ adalah himpunan koterabaikan
$\dom U=\bigcap_{i\le n}\dom X_i\in\Sigma_n$. Pada suatu titik kita perlu
memeriksa bahwa $U$ terukur terhadap $\Sigma_n$; tetapi hal ini jelas benar,
karena semua himpunan yang relevan $\{\omega:X_i(\omega)\le a\}$,
$\{\omega:X_i(\omega)\ge b\}$ termasuk dalam $\Sigma_n$. Untuk setiap
$\omega\in H$, terapkan (a) pada daftar
$X_0(\omega),\ldots,X_n(\omega)$ untuk melihat bahwa

\Centerline{$(b-a)U(\omega)\le\sum_{k=0}^{n-1}\chi F_k(\omega)
(Y_{k+1}(\omega)-Y_k(\omega))$.}

\noindent Karena $H$ koterabaikan, diperoleh

$$\eqalignno{(b-a)\Expn(U)
&\le\sum_{k=0}^{n-1}\int_{F_k}Y_{k+1}-Y_k
\le\sum_{k=0}^{n-1}\int Y_{k+1}-Y_k\cr
\noalign{\noindent (menggunakan (b-iii))}
&=\Expn(Y_n-Y_0)
\le\Expn((X_n-X_0)^+)\cr}$$

\noindent karena $Y_n-Y_0\le(X_n-X_0)^+$ di mana-mana pada
$\dom X_n\cap\dom X_0$. Ini menuntaskan pembuktian.
}%end of proof of 275F


\leader{275G}{}\cmmnt{ Sekarang kita siap untuk teorema-teorema utama dalam
bagian ini.

\medskip

\noindent}{\bf Teorema Konvergensi Martingal Doob} Misalkan
$\sequencen{X_n}$ adalah suatu martingal pada ruang probabilitas
$(\Omega,\Sigma,\mu)$, dan anggap bahwa
$\sup_{n\in\Bbb N}\Expn(|X_n|)<\infty$. Maka
$\lim_{n\to\infty}X_n(\omega)$ terdefinisi dalam $\Bbb R$ untuk hampir
setiap $\omega$ dalam $\Omega$.

\proof{{\bf (a)} Tetapkan $H=\bigcap_{n\in\Bbb N}\dom X_n$, dan untuk
$\omega\in H$ tetapkan

\Centerline{$Y(\omega)=\liminf_{n\to\infty}X_n(\omega)$,
\quad$Z(\omega)=\limsup_{n\to\infty}X_n(\omega)$,}

\noindent dengan memperbolehkan $\pm\infty$ dalam kedua kasus. Namun
perhatikan bahwa $Y\le\liminf_{n\to\infty}|X_n|$, sehingga berdasarkan lemma
Fatou $Y(\omega)<\infty$ untuk hampir setiap $\omega$; demikian pula
$Z(\omega)>-\infty$ untuk hampir setiap $\omega$. Karena itu cukup jika saya
dapat menunjukkan bahwa $Y\eae Z$, sebab dengan demikian
$Y(\omega)=Z(\omega)\in\Bbb R$ untuk hampir setiap $\omega$, dan
$\sequencen{X_n(\omega)}$ akan konvergen untuk hampir setiap $\omega$.

\medskip

{\bf (b)} \Quer\ Jadi andaikan, jika mungkin, bahwa $Y$ dan $Z$ tidak sama
hampir di mana-mana. Tentu keduanya terukur terhadap $\hat\Sigma$, dengan
$(\Omega,\hat\Sigma,\hat\mu)$ adalah pelengkapan
$(\Omega,\Sigma,\mu)$, sehingga kita harus mempunyai

\Centerline{$\hat\mu\{\omega:\omega\in H,\,Y(\omega)<Z(\omega)\}>0$.}

\noindent Oleh karena itu terdapat bilangan rasional $q$, $q'$ sedemikian
sehingga $q<q'$ dan $\hat\mu G>0$, dengan

\Centerline{$G=\{\omega:\omega\in H,\,Y(\omega)<q<q'<Z(\omega)\}$.}

\noindent Sekarang, untuk setiap $\omega\in H$ dan $n\in\Bbb N$, misalkan
$U_n(\omega)$ adalah banyaknya pelintasan naik dari $q$ ke $q'$ dalam daftar
$X_0(\omega),\ldots,
\ifdim\pagewidth>467pt\penalty-50\fi
X_n(\omega)$. Maka 275F memberi tahu kita bahwa

\Centerline{$\Expn(U_n)\le\Bover1{q'-q}\Expn((X_n-X_0)^+)
\le\Bover1{q'-q}\Expn(|X_n|+|X_0|)
\le\Bover{2M}{q'-q}$,}

\noindent jika kita menulis $M=\sup_{i\in\Bbb N}\Expn(|X_i|)$. Berdasarkan
teorema B.Levi, $U(\omega)=\sup_{n\in\Bbb N}U_n(\omega)<\infty$ untuk hampir
setiap $\omega$. Di sisi lain, jika $\omega\in G$, maka terdapat $j$, $k$
sebesar apa pun sedemikian sehingga $X_j(\omega)<q$ dan $X_k(\omega)>q'$,
sehingga $U(\omega)=\infty$. Ini berarti bahwa $\hat\mu G$ harus bernilai
$0$, bertentangan dengan pilihan $q$ dan $q'$. \Bang

\medskip

{\bf (c)} Jadi sebenarnya kita harus mempunyai $Y\eae Z$, dan
$\sequencen{X_n(\omega)}$ konvergen untuk hampir setiap $\omega$, sebagaimana
diklaim.
}%end of proof of 275G

\leader{275H}{Teorema}
Misalkan $(\Omega,\Sigma,\mu)$ adalah ruang probabilitas, dan
$\sequencen{\Sigma_n}$ barisan subaljabar-sigma dari $\Sigma$ yang tak
menurun. Misalkan $\sequencen{X_n}$ adalah martingal yang teradaptasi terhadap
$\sequencen{\Sigma_n}$. Maka pernyataan berikut ekuivalen:

\quad(i) terdapat variabel acak $X$, dengan ekspektasi berhingga, sedemikian
sehingga $X_n$ adalah suatu ekspektasi bersyarat dari $X$ pada $\Sigma_n$
untuk setiap $n$;

\quad (ii) $\{X_n:n\in\Bbb N\}$ terintegralkan secara seragam;

\quad (iii) $X_{\infty}(\omega)=\lim_{n\to\infty}X_n(\omega)$ terdefinisi
dalam $\Bbb R$ untuk hampir setiap $\omega$, dan
$\Expn(|X_{\infty}|)=\lim_{n\to\infty}\Expn(|X_n|)<\infty$.

\proof{{\bf(i)$\Rightarrow$(ii)} Berdasarkan 246D, himpunan semua ekspektasi
bersyarat dari $X$ terintegralkan secara seragam, sehingga
$\{X_n:n\in\Bbb N\}$ tentu terintegralkan secara seragam.

\medskip

{\bf (ii)$\Rightarrow$(iii)} Jika $\{X_n:n\in\Bbb N\}$ terintegralkan
secara seragam, kita tentu mempunyai
$\sup_{n\in\Bbb N}\Expn(|X_n|)<\infty$, sehingga 275G memberi tahu kita
bahwa $X_{\infty}$ terdefinisi hampir di mana-mana. Berdasarkan 246Ja,
$X_{\infty}$ terintegralkan dan
$\lim_{n\to\infty}\Expn(|X_n-X_{\infty}|)=0$. Akibatnya
$\Expn(|X_{\infty}|)=\lim_{n\to\infty}\Expn(|X_n|)<\infty$.

\medskip

{\bf (iii)$\Rightarrow$(i)} Karena
$\Expn(|X_{\infty}|)=\lim_{n\to\infty}\Expn(|X_n|)$,
$\lim_{n\to\infty}\Expn(|X_n-X_{\infty}|)=0$ (245H(a-ii)). Sekarang ambil
$n\in\Bbb N$ dan $E\in\Sigma_n$. Maka

\Centerline{$\int_EX_n=\lim_{m\to\infty}\int_EX_m=\int_EX_{\infty}$.}

\noindent Karena $E$ sebarang, $X_n$ adalah suatu ekspektasi bersyarat dari
$X_{\infty}$ pada $\Sigma_n$.
}%end of proof of 275H

\leader{275I}{Teorema} Misalkan $(\Omega,\Sigma,\mu)$ adalah ruang
probabilitas, dan $\sequencen{\Sigma_n}$ barisan subaljabar-sigma dari
$\Sigma$ yang tak menurun; tuliskan $\Sigma_{\infty}$ untuk aljabar-sigma yang
dibangkitkan oleh $\bigcup_{n\in\Bbb N}\Sigma_n$. Misalkan $X$ adalah
sebarang variabel acak bernilai real pada $\Omega$ dengan ekspektasi berhingga,
dan untuk setiap $n\in\Bbb N$ misalkan $X_n$ adalah suatu ekspektasi bersyarat
dari $X$ pada $\Sigma_n$. Maka
$X_{\infty}(\omega)=\lim_{n\to\infty}X_n(\omega)$ terdefinisi hampir di
mana-mana; $\lim_{n\to\infty}\Expn(|X_{\infty}-X_n|)=0$, dan $X_{\infty}$
adalah suatu ekspektasi bersyarat dari $X$ pada $\Sigma_{\infty}$.

\proof{ Berdasarkan 275G-275H, kita tahu bahwa $X_{\infty}$ terdefinisi
hampir di mana-mana, dan, sebagaimana dicatat dalam bukti 275H,
$\lim_{n\to\infty}\Expn(|X_{\infty}-X_n|)=0$. Untuk melihat bahwa
$X_{\infty}$ adalah suatu ekspektasi bersyarat dari $X$ pada
$\Sigma_{\infty}$, tetapkan

\Centerline{$\Cal A
=\{E:E\in\Sigma_{\infty},\,\int_EX_{\infty}=\int_EX\}$,
\quad$\Cal I=\bigcup_{n\in\Bbb N}\Sigma_n$.}

\noindent Sekarang $\Cal I$ dan $\Cal A$ memenuhi syarat-syarat Teorema
Kelas Monoton (136B). \Prf\ {\bf ($\pmb{\alpha}$)} Tentu saja
$\Omega\in\Cal I$ dan $\Cal I$ tertutup terhadap irisan berhingga, karena
$\sequencen{\Sigma_n}$ adalah barisan aljabar-sigma yang tak menurun; bahkan
$\Cal I$ merupakan subaljabar dari $\Cal P\Omega$, dan tertutup terhadap
gabungan berhingga dan komplemen. {\bf ($\pmb{\beta}$)} Jika
$E\in\Cal I$, katakanlah $E\in\Sigma_n$; maka

\Centerline{$\int_EX_{\infty}=\lim_{m\to\infty}\int_EX_m=\int_EX$,}

\noindent seperti dalam (iii)$\Rightarrow$(i) dari 275H, sehingga
$E\in\Cal A$. Jadi $\Cal I\subseteq\Cal A$.
{\bf ($\pmb{\gamma}$)} Jika $E$, $F\in\Cal A$ dan $E\subseteq F$, maka

\Centerline{$\int_{F\setminus E}X_{\infty}
=\int_FX_{\infty}-\int_EX_{\infty}
=\int_FX-\int_EX
=\int_{F\setminus E}X$,}

\noindent sehingga $F\setminus E\in\Cal A$.
{\bf($\pmb{\delta}$)} Jika $\sequencen{E_k}$ adalah barisan tak menurun
dalam $\Cal A$ dengan gabungan $E$, maka

\Centerline{$\int_{E}X_{\infty}
=\lim_{k\to\infty}\int_{E_k}X_{\infty}
=\lim_{k\to\infty}\int_{E_k}X
=\int_{E}X$,}

\noindent sehingga $E\in\Cal A$. Jadi $\Cal A$ adalah kelas Dynkin.\ \Qed

Akibatnya, berdasarkan 136B, $\Cal A$ mencakup $\Sigma_{\infty}$; yakni,
$X_{\infty}$ adalah suatu ekspektasi bersyarat dari $X$ pada
$\Sigma_{\infty}$.
}%end of proof of 275I

\cmmnt{\medskip

\noindent{\bf Catatan} Saya telah menulis
`$\lim_{n\to\infty}\Expn(|X_n-X_{\infty}|)=0$'; tetapi Anda mungkin lebih
suka mengatakan `$X^{\ssbullet}_{\infty}=\lim_{n\to\infty}X_n^{\ssbullet}$
dalam $L^1(\mu)$', seperti dalam Bab 24.

Teorema ini begitu penting sehingga Anda mungkin tertarik pada pembuktian
yang didasarkan pada 275D alih-alih 275E-275G; lihat 275Xd.
}

\leader{*275J}{}\cmmnt{ Sebagai akibat teorema ini saya memberikan hasil
penting, sejenis teorema kepadatan untuk ukuran produk.

\medskip

\noindent}{\bf Proposisi} Misalkan
$\sequencen{(\Omega_n,\Sigma_n,\mu_n)}$ adalah barisan ruang probabilitas
dengan produk $(\Omega,\Sigma,\mu)$. Misalkan $X$ adalah variabel acak
bernilai real pada $\Omega$ dengan ekspektasi berhingga. Untuk setiap
$n\in\Bbb N$ definisikan $X_n$ dengan menetapkan

\Centerline{$X_n(\pmb{\omega})
=\int X(\omega_0,\ldots,\omega_n,\xi_{n+1},\ldots)d(\xi_{n+1},\ldots)$}

\noindent di mana pun ini terdefinisi, dengan saya menulis
`$\int\ldots d(\xi_{n+1},\ldots)$' untuk menyatakan integrasi terhadap
ukuran produk $\lambda'_n$ pada $\prod_{i\ge n+1}\Omega_i$. Maka
$X(\pmb{\omega})=\lim_{n\to\infty}X_n(\pmb{\omega})$ untuk hampir setiap
$\pmb{\omega}=(\omega_0,\omega_1,\ldots)$ dalam $\Omega$, dan
$\lim_{n\to\infty}\Expn(|X-X_n|)=0$.

\proof{ Untuk setiap $n$, kita dapat mengidentifikasi $\mu$ dengan produk
$\lambda_n$ dan $\lambda'_n$, dengan $\lambda_n$ adalah ukuran produk pada
$\Omega_0\times\ldots\times\Omega_n$ (254N). Jadi 253H memberi tahu kita
bahwa $X_n$ adalah suatu ekspektasi bersyarat dari $X$ pada aljabar-sigma
$\Lambda_n=\{E\times\prod_{i>n}\Omega_i:E\in\dom\lambda_n\}$. Karena
(kembali berdasarkan 254N) kita dapat memandang $\lambda_{n+1}$ sebagai
produk $\lambda_n$ dan $\mu_{n+1}$, $\Lambda_n\subseteq\Lambda_{n+1}$ untuk
setiap $n$. Jadi 275I memberi tahu kita bahwa $\sequencen{X_n}$ konvergen
hampir di mana-mana menuju suatu ekspektasi bersyarat $X_{\infty}$ dari $X$
pada aljabar-sigma $\Lambda_{\infty}$ yang dibangkitkan oleh
$\bigcup_{n\in\Bbb N}\Lambda_n$. Sekarang
$\Lambda_{\infty}\subseteq\Sigma$ dan juga
$\Tensorhat_{n\in\Bbb N}\Sigma_n\subseteq\Lambda_{\infty}$, sehingga setiap
anggota $\Sigma$ diapit oleh dua anggota $\Lambda_{\infty}$ yang berukuran
sama (254Ff), dan $X_{\infty}$ harus sama dengan $X$ hampir di mana-mana.
Selain itu, 275I juga memberi tahu kita bahwa

\Centerline{$\lim_{n\to\infty}\Expn(|X-X_n|)
=\lim_{n\to\infty}\Expn(|X_{\infty}-X_n|)=0$,}

\noindent sebagaimana diperlukan.
}%end of proof of 275J

\leader{275K}{Martingal \dvrocolon{terbalik}}\cmmnt{ Kita mempunyai hasil
yang bersesuaian dengan 275I untuk barisan aljabar-sigma yang {\it menurun}.
Meskipun lebih jarang digunakan daripada 275G-275I, %275G 275H 275I
hasil ini mempunyai penerapan yang sangat penting.

\medskip

\noindent}{\bf Teorema} Misalkan $(\Omega,\Sigma,\mu)$ adalah ruang
probabilitas, dan $\sequencen{\Sigma_n}$ barisan subaljabar-sigma dari
$\Sigma$ yang tak menaik, dengan irisan $\Sigma_{\infty}$. Misalkan $X$
adalah sebarang variabel acak bernilai real dengan ekspektasi berhingga, dan
untuk setiap $n\in\Bbb N$ misalkan $X_n$ adalah suatu ekspektasi bersyarat
dari $X$ pada $\Sigma_n$. Maka
$X_{\infty}=\lim_{n\to\infty}X_n$ terdefinisi hampir di mana-mana dan
merupakan suatu ekspektasi bersyarat dari $X$ pada $\Sigma_{\infty}$.

\proof{{\bf (a)} Tetapkan $H=\bigcap_{n\in\Bbb N}\dom X_n$, sehingga $H$
koterabaikan. Untuk $n\in\Bbb N$, $a<b$ dalam $\Bbb R$, dan $\omega\in H$,
tuliskan $U_{abn}(\omega)$ untuk banyaknya pelintasan naik dari $a$ ke $b$
dalam daftar $X_n(\omega),X_{n-1}(\omega),\ldots,X_0(\omega)$ (275E).
Maka

$$\eqalignno{\Expn(U_{abn})
&\le\Bover1{b-a}\Expn((X_0-X_n)^+)\cr
\displaycause{275F}
&\le\Bover1{b-a}\Expn(|X_0|+|X_n|)
\le\Bover2{b-a}\Expn(|X_0|)<\infty.\cr}$$

\noindent Jadi $\lim_{n\to\infty}U_{abn}(\omega)$ berhingga untuk hampir
setiap $\omega$. Namun ini berarti bahwa

\Centerline{$\{\omega:\liminf_{n\to\infty}X_n(\omega)<a,\,
  \limsup_{n\to\infty}X_n(\omega)>b\}$}

\noindent terabaikan. Karena $a$ dan $b$ sebarang, $\sequencen{X_n}$
konvergen hampir di mana-mana, tepat seperti dalam 275G. Tetapkan
$X_{\infty}(\omega)=\lim_{n\to\infty}X_n(\omega)$ kapan pun ini terdefinisi
dalam $\Bbb R$.

\medskip

{\bf (b)} Berdasarkan 246D, $\{X_n:n\in\Bbb N\}$ terintegralkan secara
seragam, sehingga $\Expn(|X_n-X_{\infty}|)\to 0$ ketika $n\to\infty$
(246Ja), dan

\Centerline{$\int_EX_{\infty}=\lim_{n\to\infty}\int_EX_n=\int_EX_0$}

\noindent untuk setiap $E\in\Sigma_{\infty}$.

\medskip

{\bf (c)} Sekarang terdapat himpunan koterabaikan
$G\in\Sigma_{\infty}$ sedemikian sehingga $G\subseteq\dom X_{\infty}$ dan
$X_{\infty}\restr G$ terukur terhadap $\Sigma_{\infty}$. \Prf\ Untuk setiap
$n\in\Bbb N$, terdapat himpunan koterabaikan $G_n\in\Sigma_n$ sedemikian
sehingga $G_n\subseteq\dom X_n$ dan $X_n\restr G_n$ terukur terhadap
$\Sigma_n$. Tetapkan $G'=\bigcup_{n\in\Bbb N}\bigcap_{m\ge n}G_m$; maka,
untuk setiap $r\in\Bbb N$,
$G'=\bigcup_{n\ge r}\bigcap_{m\ge n}G_m$ termasuk dalam $\Sigma_r$,
sehingga $G'\in\Sigma_{\infty}$, sedangkan tentu saja $G'$ koterabaikan.
Untuk $n\in\Bbb N$, tetapkan $X'_n(\omega)=X_n(\omega)$ untuk
$\omega\in G_n$, dan $0$ untuk $\omega\in\Omega\setminus G_n$; maka untuk
$\omega\in G'$,
$\lim_{n\to\infty}X'_n(\omega)=\lim_{n\to\infty}X_n(\omega)$ jika salah
satunya terdefinisi dalam $\Bbb R$. Dengan menulis
$X'_{\infty}=\lim_{n\to\infty}X'_n$ kapan pun ini terdefinisi dalam
$\Bbb R$, 121F dan 121H memberi tahu kita bahwa $X'_{\infty}$ terukur
terhadap $\Sigma_r$ dan $\dom X'_{\infty}\in\Sigma_r$ untuk setiap
$r\in\Bbb N$, sehingga $G''=\dom X'_{\infty}$ termasuk dalam
$\Sigma_{\infty}$ dan $X'_{\infty}$ terukur terhadap $\Sigma_{\infty}$.
Kita juga tahu, dari (a), bahwa $G''$ koterabaikan. Jadi dengan menetapkan
$G=G'\cap G''$ kita memperoleh hasilnya.\ \Qed

Dengan demikian $X_{\infty}$ adalah suatu ekspektasi bersyarat dari $X$ pada
$\Sigma_{\infty}$.
}%end of proof of 275K

\leader{275L}{Waktu\dvro{ henti: }{ henti}}\cmmnt{ Dalam suatu arti,
pekerjaan utama bagian ini telah selesai; saya tidak mempunyai ruang bagi
teorema penting lain yang sebanding dengan 275G-275I. Namun tidak tepat untuk
meninggalkan bab ini tanpa mendeskripsikan secara singkat salah satu gagasan
paling subur dalam bidang ini.

\medskip

\noindent}{\bf Definisi} Misalkan $(\Omega,\Sigma,\mu)$ adalah ruang
probabilitas, dengan pelengkapan $(\Omega,\hat\Sigma,\hat\mu)$, dan
$\sequencen{\Sigma_n}$ barisan subaljabar-sigma dari $\hat\Sigma$ yang tak
menurun. Suatu {\bf waktu henti yang teradaptasi terhadap
$\sequencen{\Sigma_n}$}\cmmnt{ (juga disebut `{\bf waktu opsional}',
`{\bf waktu Markov}')} adalah fungsi $\tau$ dari $\Omega$ ke
$\Bbb N\cup\{\infty\}$ sedemikian sehingga
$\{\omega:\tau(\omega)\le n\}\in\Sigma_n$ untuk setiap $n\in\Bbb N$.

\cmmnt{\medskip

\noindent{\bf Catatan} Tentu saja syarat

\Centerline{$\{\omega:\tau(\omega)\le n\}\in\Sigma_n$ untuk setiap
$n\in\Bbb N$}

\noindent dapat digantikan oleh syarat yang ekuivalen

\Centerline{$\{\omega:\tau(\omega)=n\}\in\Sigma_n$
untuk setiap $n\in\Bbb N$.}

\noindent Saya mendahulukan ungkapan pertama karena lebih sesuai untuk
himpunan indeks lain (lihat 275Cc).
}%end of comment

\leader{275M}{Contoh (a)} Jika $\sequencen{X_n}$ adalah martingal yang
teradaptasi terhadap barisan aljabar-sigma $\sequencen{\Sigma_n}$, dan
${H_n}$ adalah himpunan bagian Borel dari $\BbbR^{n+1}$ untuk setiap $n$,
maka kita mempunyai waktu henti $\tau$ yang teradaptasi terhadap
$\sequencen{\Sigma_n}$ dan didefinisikan oleh rumus

\Centerline{$\tau(\omega)
=\inf\{n:\omega\in\bigcap_{i\le n}\dom X_i,\,
  (X_0(\omega),\ldots,X_n(\omega))\in H_n\}$,}

\noindent dengan menetapkan $\inf\emptyset=\infty$\cmmnt{ seperti biasa}.
\cmmnt{(Sebab berdasarkan 121Ka, himpunan
$E_n=\{\omega:(X_0(\omega),\ldots,X_n(\omega))\in H_n\}$ termasuk dalam
$\Sigma_n$ untuk setiap $n$, dan
$\{\omega:\tau(\omega)\le n\}=\bigcup_{i\le n}E_i$.)} Khususnya,
\cmmnt{misalnya,} rumus-rumus

\Centerline{$\inf\{n:X_n(\omega)\ge a\}$, \quad
$\inf\{n:|X_n(\omega)|> a\}$}

\noindent mendefinisikan waktu henti.

\header{275Mb}{\bf (b)} Setiap fungsi konstan
$\tau:\Omega\to\Bbb N\cup\{\infty\}$ adalah waktu henti. Jika $\tau$,
$\tau'$ adalah dua waktu henti yang teradaptasi terhadap barisan aljabar-sigma
$\sequencen{\Sigma_n}$ yang sama, maka $\tau\wedge\tau'$ adalah waktu henti
yang teradaptasi terhadap $\sequencen{\Sigma_n}$, dengan menetapkan
$(\tau\wedge\tau')(\omega)=\min(\tau(\omega),\tau'(\omega))$ untuk
$\omega\in\Omega$\prooflet{, karena

\Centerline{$\{\omega:(\tau\wedge\tau')(\omega)\le n\}
=\{\omega:\tau(\omega)\le n\}\cup\{\omega:\tau'(\omega)\le n\}
\in\Sigma_n$}

\noindent untuk setiap $n\in\Bbb N$}.

\leader{275N}{Lemma} Misalkan $(\Omega,\Sigma,\mu)$ adalah ruang
probabilitas lengkap, dan $\sequencen{\Sigma_n}$ barisan subaljabar-sigma dari
$\Sigma$ yang tak menurun. Anggap bahwa $\tau$ dan $\tau'$ adalah waktu henti
pada $\Omega$, dan $\sequencen{X_n}$ suatu martingal, semuanya teradaptasi
terhadap $\sequencen{\Sigma_n}$.

(a) Keluarga

\Centerline{$\tilde\Sigma_{\tau}
=\{E:E\in\Sigma,\,E\cap\{\omega:\tau(\omega)\le n\}\in\Sigma_n$ untuk
setiap $n\in\Bbb N\}$}

\noindent adalah subaljabar-sigma dari $\Sigma$.

(b) Jika $\tau(\omega)\le\tau'(\omega)$ untuk setiap $\omega$, maka
$\tilde\Sigma_{\tau}\subseteq\tilde\Sigma_{\tau'}$.

(c) Sekarang anggap bahwa $\tau$ berhingga hampir di mana-mana. Tetapkan

\Centerline{$\tilde X_{\tau}(\omega)=X_{\tau(\omega)}(\omega)$}

\noindent kapan pun $\tau(\omega)<\infty$ dan
$\omega\in\dom X_{\tau(\omega)}$. Maka
$\dom\tilde X_{\tau}\in\tilde\Sigma_{\tau}$ dan $\tilde X_{\tau}$ terukur
terhadap $\tilde\Sigma_{\tau}$.

(d) Jika $\tau$ terbatas secara esensial, yakni, terdapat suatu
$m\in\Bbb N$ sedemikian sehingga $\tau\le m$ hampir di mana-mana, maka
$\Expn(\tilde X_{\tau})$ ada dan sama dengan $\Expn(X_0)$.

(e) Jika $\tau\le\tau'$ hampir di mana-mana, dan $\tau'$ terbatas secara
esensial, maka $\tilde X_{\tau}$ adalah suatu ekspektasi bersyarat dari
$\tilde X_{\tau'}$ pada $\tilde\Sigma_{\tau}$.

\proof{{\bf (a)} Hal ini elementer. Tuliskan
$H_n=\{\omega:\tau(\omega)\le n\}$ untuk setiap $n\in\Bbb N$. Himpunan kosong
termasuk dalam $\tilde\Sigma_{\tau}$ karena termasuk dalam $\Sigma_n$ untuk
setiap $n$.
Jika $E\in\tilde\Sigma_{\tau}$, maka

\Centerline{$(\Omega\setminus E)\cap H_n=H_n\setminus(E\cap
H_n)\in\Sigma_n$}

\noindent karena $H_n\in\Sigma_n$; ini benar untuk setiap $n$, sehingga
$\Omega\setminus E\in\tilde\Sigma_{\tau}$. Jika $\sequence{k}{E_k}$ adalah
sebarang barisan dalam $\tilde\Sigma_{\tau}$, maka

\Centerline{$(\bigcup_{k\in\Bbb N}E_k)\cap H_n=\bigcup_{k\in\Bbb
N}E_k\cap H_n\in\Sigma_n$}

\noindent untuk setiap $n$, sehingga
$\bigcup_{k\in\Bbb N}E_k\in\tilde\Sigma_{\tau}$.

\medskip

{\bf (b)} Jika $E\in\tilde\Sigma_{\tau}$ maka tentu saja $E\in\Sigma$,
dan jika $n\in\Bbb N$ maka
$\{\omega:\tau'(\omega)\le n\}\subseteq\{\omega:\tau(\omega)\le n\}$,
sehingga

\Centerline{$E\cap\{\omega:\tau'(\omega)\le n\}
=E\cap\{\omega:\tau(\omega)\le n\}\cap\{\omega:\tau'(\omega)\le n\}$}

\noindent termasuk dalam $\Sigma_n$; karena $n$ sebarang,
$E\in\tilde\Sigma_{\tau'}$.

\medskip

{\bf (c)} Tetapkan $H_n=\{\omega:\tau(\omega)\le n\}$ untuk setiap
$n\in\Bbb N$. Untuk setiap $a\in\Bbb R$,

$$\eqalign{H_n\cap\{\omega:\omega\in
&\dom\tilde X_{\tau},\,\tilde X_{\tau}(\omega)\le a\}\cr
&=\bigcup_{k\le n}\{\omega:\tau(\omega)=k,\,\omega\in\dom X_k,
  \,X_k(\omega)\le a\}
\in\Sigma_n.\cr}$$

\noindent Karena $n$ sebarang,

\Centerline{$G_a=\{\omega:\omega\in\dom\tilde X_{\tau},\,
\tilde X_{\tau}(\omega)\le a\}\in\tilde\Sigma_{\tau}$.}

\noindent Karena $a$ sebarang,
$\dom\tilde X_{\tau}=\bigcup_{m\in\Bbb N}G_m\in\tilde\Sigma_{\tau}$ dan
$\tilde X_{\tau}$ terukur terhadap $\tilde\Sigma_{\tau}$.

\medskip

{\bf (d)} Tetapkan $H_k=\{\omega:\tau(\omega)=k\}$ untuk $k\le m$. Maka
$\bigcup_{k\le m}H_k$ koterabaikan, sehingga

\Centerline{$\Expn(\tilde X_{\tau})
=\sum_{k=0}^m\int_{H_k}X_k
=\sum_{k=0}^m\int_{H_k}X_m
=\int_{\Omega}X_m
=\int_{\Omega}X_0$.}

\medskip

{\bf (e)} Anggap bahwa $\tau'\le n$ hampir di mana-mana. Tetapkan
$H_k=\{\omega:\tau(\omega)=k\}$,
$H'_k=\{\omega:\tau'(\omega)=k\}$ untuk setiap $k$; maka baik
$\langle H_k\rangle_{k\le n}$ maupun $\langle H'_k\rangle_{k\le n}$ adalah
partisi dari himpunan bagian koterabaikan dari $\Omega$.
Sekarang anggap bahwa $E\in\tilde\Sigma_{\tau}$. Maka

\Centerline{$\int_E\tilde X_{\tau}
=\sum_{k=0}^n\int_{E\cap H_k}\tilde X_{\tau}
=\sum_{k=0}^n\int_{E\cap H_k}X_k
=\sum_{k=0}^n\int_{E\cap H_k}X_n
=\int_EX_n$}

\noindent karena $E\cap H_k\in\Sigma_k$ untuk setiap $k$. Berdasarkan (b),
$E\in\tilde\Sigma_{\tau'}$, sehingga kita juga mempunyai
$\int_E\tilde X_{\tau'}=\int_EX_n$. Jadi
$\int_E\tilde X_{\tau}=\int_E\tilde X_{\tau'}$ untuk setiap
$E\in\tilde\Sigma_{\tau}$, sebagaimana diklaim.
}%end of proof of 275N

\leader{275O}{Proposisi} Misalkan $\sequencen{X_n}$ adalah martingal dan
$\tau$ suatu waktu henti, keduanya teradaptasi terhadap barisan aljabar-sigma
$\sequencen{\Sigma_n}$ yang sama. Untuk setiap $n$, tetapkan
$(\tau\wedge n)(\omega)=\min(\tau(\omega),n)$ untuk $\omega\in\Omega$; maka
$\tau\wedge n$ adalah waktu henti, dan $\sequencen{\tilde X_{\tau\wedge n}}$
adalah martingal yang teradaptasi terhadap
$\sequencen{\tilde\Sigma_{\tau\wedge n}}$, dengan
$\tilde X_{\tau\wedge n}$ dan $\tilde\Sigma_{\tau\wedge n}$ didefinisikan
seperti dalam 275N.

\proof{ Sebagaimana dicatat dalam 275Mb, setiap $\tau\wedge n$ adalah waktu
henti. Jika $m\le n$, maka
$\tilde\Sigma_{\tau\wedge m}\subseteq\tilde\Sigma_{\tau\wedge n}$
berdasarkan 275Nb. Setiap $\tilde X_{\tau\wedge m}$ terukur terhadap
$\tilde\Sigma_{\tau\wedge m}$, dengan domain termasuk dalam
$\tilde\Sigma_{\tau\wedge m}$, berdasarkan 275Nc, dan mempunyai ekspektasi
berhingga, berdasarkan 275Nd; akhirnya, jika $m\le n$, maka
$\tilde X_{\tau\wedge m}$ adalah suatu ekspektasi bersyarat dari
$\tilde X_{\tau\wedge n}$ pada $\tilde\Sigma_{\tau\wedge m}$, berdasarkan
275Ne.
}%end of proof of 275O

\leader{275P}{Akibat} Anggap bahwa $(\Omega,\Sigma,\mu)$ adalah ruang
probabilitas dan $\sequencen{X_n}$ suatu martingal pada $\Omega$ sedemikian
sehingga $W=\sup_{n\in\Bbb N}|X_{n+1}-X_n|$ berhingga hampir di mana-mana dan
mempunyai ekspektasi berhingga. Maka untuk hampir setiap $\omega\in\Omega$,
entah $\lim_{n\to\infty}X_n(\omega)$ ada dalam $\Bbb R$, atau
$\sup_{n\in\Bbb N}X_n(\omega)=\infty$ dan
$\inf_{n\in\Bbb N}X_n(\omega)=-\infty$.

\proof{ Misalkan $\sequencen{\Sigma_n}$ adalah barisan aljabar-sigma tak
menurun yang terhadapnya $\sequencen{X_n}$ teradaptasi. Misalkan $H$ adalah
himpunan koterabaikan
$\bigcap_{n\in\Bbb N}\dom X_n\cap\{\omega:W(\omega)<\infty\}$. Untuk setiap
$m\in\Bbb N$, tetapkan

\Centerline{$\tau_m(\omega)=\inf\{n:\omega\in\dom
X_n,\,X_n(\omega)>m\}$.}

\noindent Seperti dalam 275Ma, $\tau_m$ adalah waktu henti yang teradaptasi
terhadap $\sequencen{\Sigma_n}$. Tetapkan

\Centerline{$Y_{mn}=\tilde X_{\tau_m\wedge n}$,}

\noindent yang didefinisikan seperti dalam 275O, sehingga
$\sequencen{Y_{mn}}$ adalah martingal. Jika $\omega\in H$, maka entah
$\tau_m(\omega)>n$ dan

\Centerline{$Y_{mn}(\omega)=X_n(\omega)\le m$,}

\noindent atau $0<\tau_m(\omega)\le n$ dan

\Centerline{$Y_{mn}(\omega)=X_{\tau_m(\omega)}(\omega)
\le W(\omega)+X_{\tau_m(\omega)-1}(\omega)\le W(\omega)+m$,}

\noindent atau $\tau_m(\omega)=0$ dan

\Centerline{$Y_{mn}(\omega)=X_0(\omega)$.}

\noindent Jadi

\Centerline{$Y_{mn}(\omega)\le|X_0(\omega)|+W(\omega)+m$}

\noindent untuk setiap $\omega\in H$, dan

\Centerline{$|Y_{mn}(\omega)|=2\max(0,Y_{mn}(\omega))-Y_{mn}(\omega)
\le 2(|X_0(\omega)|+W(\omega)+m)-Y_{mn}(\omega)$,}

\Centerline{$\Expn(|Y_{mn}|)\le 2\Expn(|X_0|)+2\Expn(W)+2m-\Expn(Y_{mn})
= 2\Expn(|X_0|)+2\Expn(W)+2m-\Expn(X_0)$}

\noindent berdasarkan 275Nd. Karena ini benar untuk setiap $n\in\Bbb N$,
$\sup_{n\in\Bbb N}\Expn(|Y_{mn}|)<\infty$, dan
$\lim_{n\to\infty}Y_{mn}$ terdefinisi dalam $\Bbb R$ hampir di mana-mana,
berdasarkan Teorema Konvergensi Martingal Doob (275G). Misalkan $F_m$ adalah
himpunan koterabaikan tempat $\sequencen{Y_{mn}}$ konvergen.
Tetapkan $H^*=H\cap\bigcap_{m\in\Bbb N}F_m$, sehingga $H^*$ koterabaikan.

Sekarang tinjau

\Centerline{$E=\{\omega:\omega\in H^*,\,\sup_{n\in\Bbb
N}X_n(\omega)<\infty\}$.}

\noindent Untuk setiap $\omega\in E$, harus terdapat $m\in\Bbb N$ sedemikian
sehingga $\sup_{n\in\Bbb N}X_n(\omega)\le m$. Sekarang ini berarti bahwa
$Y_{mn}(\omega)=X_n(\omega)$ untuk setiap $n$, dan karena $\omega\in F_m$
kita mempunyai

\Centerline{$\lim_{n\to\infty}X_n(\omega)
=\lim_{n\to\infty}Y_{mn}(\omega)\in\Bbb R$.}

\noindent Ini berarti bahwa $\sequencen{X_n(\omega)}$ konvergen untuk hampir
setiap $\omega$ sedemikian sehingga $\{X_n(\omega):n\in\Bbb N\}$ terbatas
ke atas.

Dengan cara serupa, $\sequencen{X_n(\omega)}$ konvergen untuk hampir setiap
$\omega$ sedemikian sehingga $\{X_n(\omega):n\in\Bbb N\}$ terbatas ke bawah,
yang menuntaskan pembuktian.
}%end of proof of 275P


\exercises{
\leader{275X}{Latihan dasar $\pmb{>}$(a)}
%\sqheader 275Xa
Misalkan $\sequencen{X_n}$ adalah barisan variabel acak yang saling bebas
dengan ekspektasi nol dan varians berhingga. Tetapkan
$s_n=(\sum_{i=0}^n\Var(X_i))^{1/2}$,
$Y_n=(X_0+\ldots+X_n)^2-s_n^2$ untuk setiap $n$. Tunjukkan bahwa
$\sequencen{Y_n}$ adalah martingal.
%275B

\sqheader 275Xb Misalkan $\sequencen{X_n}$ adalah suatu martingal. Tunjukkan
bahwa untuk setiap $\epsilon>0$,
$\Pr(\sup_{n\in\Bbb N}|X_n|\ge\epsilon)
\le\bover1\epsilon\sup_{n\in\Bbb N}\Expn(|X_n|)$.
%275D

\spheader 275Xc {\bf Skema guci P\'olya} Bayangkan sebuah kotak yang memuat
bola merah dan putih. Pada setiap langkah, sebuah bola diambil secara acak
dari kotak lalu dikembalikan bersama sebuah bola lain dengan warna yang sama.
(i) Dengan $R_n$, $W_n$ menyatakan banyaknya bola merah dan putih setelah
langkah ke-$n$, dan $X_n=R_n/(R_n+W_n)$, tunjukkan bahwa
$\sequencen{X_n}$ adalah martingal. (ii) Dengan memulai dari $R_0=W_0=1$,
temukan distribusi $(R_n,W_n)$ untuk setiap $n$. (iii) Tunjukkan bahwa
$X=\lim_{n\to\infty}X_n$ terdefinisi hampir di mana-mana, dan temukan
distribusinya ketika $R_0=W_0=1$. (Lihat {\smc Feller 66} untuk pembahasan
nilai awal lainnya.)
%275G

\sqheader 275Xd Misalkan $(\Omega,\Sigma,\mu)$ adalah ruang probabilitas,
dan $\sequencen{\Sigma_n}$ barisan subaljabar-sigma dari $\Sigma$ yang tak
menurun; untuk setiap $n\in\Bbb N$ misalkan $P_n:L^1\to L^1$ adalah operator
ekspektasi bersyarat yang bersesuaian dengan $\Sigma_n$, dengan
$L^1=L^1(\mu)$ (242J). (i) Tunjukkan bahwa
$V=\{u:u\in L^1,\,\lim_{n\to\infty}\|P_nu-u\|_1=0\}$ adalah subruang linear
tertutup-$\|\,\|_1$ dari $L^1$. (ii) Tunjukkan bahwa
$\{E:E\in\Sigma,\,\chi E^{\ssbullet}\in V\}$ adalah kelas Dynkin yang
mencakup $\bigcup_{n\in\Bbb N}\Sigma_n$, sehingga mencakup aljabar-sigma
$\Sigma_{\infty}$ yang dibangkitkan oleh
$\bigcup_{n\in\Bbb N}\Sigma_n$. (iii) Tunjukkan bahwa jika $u\in L^1$ maka
$v=\sup_{n\in\Bbb N}P_n|u|$ terdefinisi dalam $L^1$ dan berbentuk
$W^{\ssbullet}$ dengan $\Pr(W\ge\epsilon)\le\bover1{\epsilon}\|u\|_1$ untuk
setiap $\epsilon>0$. \Hint{275D.} (iv) Tunjukkan bahwa jika $X$ adalah
variabel acak terukur terhadap $\Sigma_{\infty}$ dengan ekspektasi berhingga,
dan untuk setiap $n\in\Bbb N,\ X_n$ adalah suatu ekspektasi bersyarat dari
$X$ pada $\Sigma_n$, maka $X^{\ssbullet}\in V$ dan
$X\eae\lim_{n\to\infty}X_n$.
\Hint{terapkan (iii) pada $u=(X-X_m)^{\ssbullet}$ untuk $m$ yang besar.}
%275I

\spheader 275Xe Misalkan $(\Omega,\Sigma,\mu)$ adalah ruang probabilitas,
$\sequencen{\Sigma_n}$ barisan subaljabar-sigma dari $\Sigma$ yang tak
menurun, dan $\Sigma_{\infty}$ aljabar-sigma yang dibangkitkan oleh
$\bigcup_{n\in\Bbb N}\Sigma_n$. Untuk setiap
$n\in\Bbb N\cup\{\infty\}$ misalkan $P_n:L^1\to L^1$ adalah operator
ekspektasi bersyarat yang bersesuaian dengan $\Sigma_n$, dengan
$L^1=L^1(\mu)$. Tunjukkan bahwa
$\lim_{n\to\infty}\|P_nu-P_{\infty}u\|_p=0$ kapan pun
$p\in\coint{1,\infty}$ dan $u\in L^p(\mu)$.
\Hint{275Xd, 233J/242K, 246Xg.}
%275Xd, 275I

\spheader 275Xf Misalkan $\sequencen{X_n}$ adalah suatu martingal, dan
anggap bahwa $p\in\ooint{1,\infty}$ sedemikian sehingga
$\sup_{n\in\Bbb N}\|X_n\|_p<\infty$. Tunjukkan bahwa
$X=\lim_{n\to\infty}X_n$ terdefinisi hampir di mana-mana dan bahwa
$\lim_{n\to\infty}\|X_n-X\|_p=0$.
%275Xe, 275H, 275I

\sqheader 275Xg Misalkan $(\Omega,\Sigma,\mu)$ adalah $[0,1]$ dengan ukuran
Lebesgue. Untuk setiap $n\in\Bbb N$ misalkan $\Sigma_n$ adalah subaljabar
berhingga dari $\Sigma$ yang dibangkitkan oleh interval jenis $[0,2^{-n}r]$
untuk $r\le 2^n$. Gunakan 275I untuk menunjukkan bahwa untuk setiap
$X:[0,1]\to\Bbb R$ yang terintegralkan kita harus mempunyai
$X(t)=\lim_{n\to\infty}2^n\int_{I_n(t)}X$ untuk hampir setiap
$t\in\coint{0,1}$, dengan $I_n(t)$ adalah interval berbentuk
$\coint{2^{-n}r,2^{-n}(r+1)}$ yang memuat $t$. Bandingkan hasil ini dengan
223A dan 261Yd.
%275I

\spheader 275Xh Dalam 275K, tunjukkan bahwa
$\lim_{n\to\infty}\|X_n-X_{\infty}\|_p=0$ untuk setiap
$p\in\coint{1,\infty}$ sedemikian sehingga $\|X_0\|_p$ berhingga.
(Bandingkan 275Xe.)
%275K

\spheader 275Xi\dvAnew{2012} Misalkan $(\Omega,\Sigma,\mu)$ adalah ruang
probabilitas, dengan pelengkapan $(\Omega,\hat\Sigma,\hat\mu)$, dan
$\sequencen{\Sigma_n}$ barisan subaljabar-sigma dari $\hat\Sigma$ yang tak
menurun. Tunjukkan bahwa jika $\sequence{i}{\tau_i}$ adalah barisan waktu
henti yang teradaptasi terhadap $\sequencen{\Sigma_n}$, dan kita menetapkan
$\tau(\omega)=\sup_{i\in\Bbb N}\tau_i(\omega)$ untuk $\omega\in\Omega$,
maka $\tau$ adalah waktu henti yang teradaptasi terhadap
$\sequencen{\Sigma_n}$.
%275M

\spheader 275Xj Misalkan $(\Omega,\Sigma,\mu)$ adalah ruang probabilitas,
dengan pelengkapan $(\Omega,\hat\Sigma,\hat\mu)$, dan
$\sequencen{\Sigma_n}$ barisan subaljabar-sigma dari $\hat\Sigma$ yang tak
menurun. Misalkan $\sequencen{X_n}$ adalah martingal terintegralkan secara
seragam yang teradaptasi terhadap $\sequencen{\Sigma_n}$, dan tetapkan
$X_{\infty}=\lim_{n\to\infty}X_n$. Misalkan $\tau$ adalah waktu henti yang
teradaptasi terhadap $\sequencen{\Sigma_n}$, dan tetapkan
$\tilde X_{\tau}(\omega)=X_{\tau(\omega)}(\omega)$ kapan pun
$\omega\in\dom X_{\tau(\omega)}$, dengan memperbolehkan $\infty$ sebagai
nilai $\tau(\omega)$. Tunjukkan bahwa $\tilde X_{\tau}$ adalah suatu
ekspektasi bersyarat dari $X_{\infty}$ pada $\tilde\Sigma_{\tau}$,
sebagaimana didefinisikan dalam 275N.
%275N

\spheader 275Xk Misalkan $(\Omega,\Sigma,\mu)$ adalah ruang probabilitas,
dengan pelengkapan $(\Omega,\hat\Sigma,\hat\mu)$, dan
$\sequencen{\Sigma_n}$ barisan subaljabar-sigma dari $\hat\Sigma$ yang tak
menurun. Misalkan $\sequencen{X_n}$ adalah martingal dan $\tau$ suatu waktu
henti, keduanya teradaptasi terhadap $\sequencen{\Sigma_n}$. Anggap bahwa
$\sup_{n\in\Bbb N}\Expn(|X_n|)<\infty$ dan bahwa $\tau$ berhingga hampir di
mana-mana. Tunjukkan bahwa $\tilde X_{\tau}$, sebagaimana didefinisikan dalam
275Nc, mempunyai ekspektasi berhingga, tetapi
$\Expn(\tilde X_{\tau})$ tidak harus sama dengan $\Expn(X_0)$.
%275N

\spheader 275Xl(i)\dvAnew{2014} Tunjukkan bahwa jika $\sequencen{X_n}$ adalah
martingal sedemikian sehingga $\sup_{n\in\Bbb N}\Expn(|X_n|)$ berhingga dan
$\sequencen{X_n}$ konvergen dalam ukuran, maka $\sequencen{X_n}$ konvergen
hampir di mana-mana. (ii) Temukan martingal $\sequencen{X_n}$ sedemikian
sehingga $\sequencen{X_{2n}}\to 0$ hampir di mana-mana tetapi
$|X_{2n+1}|\ge 1$ hampir di mana-mana untuk setiap $n\in\Bbb N$.
(iii) Temukan martingal yang konvergen dalam ukuran tetapi tidak konvergen
hampir di mana-mana.
%275D+  

\leader{275Y}{Latihan lanjutan (a)}%
%\spheader 275Ya
\dvAnew{2012} Misalkan $(\Omega,\Sigma,\mu)$ adalah ruang probabilitas,
$\sequencen{\Sigma_n}$ barisan aljabar-sigma yang saling bebas dari $\Sigma$,
dan $X$ variabel acak pada $\Omega$ dengan varians berhingga. Misalkan $X_n$
adalah suatu ekspektasi bersyarat dari $X$ pada $\Sigma_n$ untuk setiap $n$.
Tunjukkan bahwa $\lim_{n\to\infty}X_n=\Expn(X)$ hampir di mana-mana.
\Hint{tinjau $\sum_{n=0}^{\infty}\Var(X_n)$.}
%-

\spheader 275Yb
Misalkan $(\Omega,\Sigma,\mu)$ adalah ruang probabilitas lengkap,
$\sequencen{\Sigma_n}$ barisan subaljabar-sigma dari $\Sigma$ yang tak
menurun dan semuanya memuat setiap himpunan terabaikan, dan
$\sequencen{X_n}$ suatu martingal yang teradaptasi terhadap
$\sequencen{\Sigma_n}$. Misalkan $\nu$ adalah ukuran probabilitas lain dengan
domain $\Sigma$ yang kontinu mutlak terhadap $\mu$, dengan turunan
Radon-Nikod\'ym $Z$. Untuk setiap $n\in\Bbb N$ misalkan $Z_n$ adalah suatu
ekspektasi bersyarat dari $Z$ pada $\Sigma_n$ (terhadap ukuran $\mu$).
(i) Tunjukkan bahwa $Z_n$ adalah turunan Radon-Nikod\'ym dari
$\nu\restr\Sigma_n$ terhadap $\mu\restr\Sigma_n$, untuk setiap $n\in\Bbb N$.
(ii) Dengan mendefinisikan $X_n/Z_n$ seperti dalam 121E, sehingga domainnya
adalah $\{\omega:\omega\in\dom X_n\cap\dom Z_n$, $Z_n(\omega)\ne 0\}$,
tunjukkan bahwa $\sequencen{X_n/Z_n}$ adalah martingal terhadap ukuran $\nu$.
%275B

\spheader 275Yc Gabungkan gagasan 275Cc dengan gagasan 275Cd-275Ce untuk
mendeskripsikan pengertian `martingal yang diindeks oleh $I$', dengan $I$
sebarang himpunan terurut parsial.
%275C

\spheader 275Yd\dvAformerly{2{}75Yc}
Misalkan $\sequence{k}{X_k}$ adalah martingal pada ruang probabilitas lengkap
$(\Omega,\Sigma,\mu)$, dan tetapkan $n\in\Bbb N$. Tetapkan
$X^*=\max(|X_0|,\ldots,\discretionary{}{}{}|X_n|)$. Misalkan
$p\in\ooint{1,\infty}$. Tunjukkan bahwa
$\|X^*\|_p\le\Bover{p}{p-1}\|X_n\|_p$. ({\it Petunjuk\/}: tetapkan
$F_t=\{\omega:X^*(\omega)\ge t\}$. Tunjukkan bahwa
$t\mu F_t\le\int_{F_t}|X_n|$. Dengan menggunakan teorema Fubini pada
$\Omega\times\coint{0,\infty}$ dan pada
$\Omega\times\coint{0,\infty}\times\coint{0,\infty}$, tunjukkan bahwa

\Centerline{$\Expn((X^*)^p)=p\int_0^{\infty}t^{p-1}\mu F_tdt$,}

\Centerline{$\int_0^{\infty}t^{p-2}\int_{F_t}|X_n|dt
=\Bover1{p-1}\Expn(|X_n|\times(X^*)^{p-1})$,}

\Centerline{$\Expn(|X_n|\times(X^*)^{p-1})
\le\|X_n\|_p\|X^*\|_p^{p-1}$.}

\noindent Bandingkan 286A di bawah.)
%275D

\spheader 275Ye\dvArevised{2012}(i)
Tunjukkan bahwa jika $a$, $b\ge 0$ maka
$a\ln^+b\le a\ln^+a+\Bover{b}e$, dengan $\ln^+t=0$ jika $t\le 1$,
dan $\ln t$ jika $t\ge 1$.
(ii) Misalkan $(\Omega,\Sigma,\mu)$ adalah ruang probabilitas lengkap dan
$X$, $Y$ variabel acak nonnegatif pada $\Omega$ sedemikian sehingga
$t\mu F_t\le\int_{F_t}X$ untuk setiap $t\ge 0$, dengan
$F_t=\{\omega:Y(\omega)\ge t\}$. Tunjukkan bahwa
$\int_{F_1}Y\le\int_{F_1}X\times\ln^+Y$, dan karenanya
$\Expn(Y)\le\Bover{e}{e-1}(1+\Expn(X\times\ln^+X))$. (iii) Tunjukkan bahwa
jika $\sequencen{X_n}$ adalah suatu martingal pada $\Omega$, $n\in\Bbb N$
dan $X^*=\sup_{i\le n}|X_i|$, maka
$\Expn(X^*)\le\Bover{e}{e-1}(1+\Expn(|X_n|\times\ln^+|X_n|))$.
%275Yd, 275D {\smc Doob 53}

\spheader 275Yf Misalkan $(\Omega,\Sigma,\mu)$ adalah ruang probabilitas
dan $\langle\Sigma_i\rangle_{i\in I}$ keluarga terhitung subaljabar-sigma
dari $\Sigma$ sedemikian sehingga untuk setiap $i$, $j\in I$, entah
$\Sigma_i\subseteq\Sigma_j$ atau $\Sigma_j\subseteq\Sigma_i$. Misalkan $X$
adalah variabel acak bernilai real pada $\Omega$ sedemikian sehingga
$\|X\|_p<\infty$, dengan $1<p<\infty$, dan anggap bahwa $X_i$ adalah suatu
ekspektasi bersyarat dari $X$ pada $\Sigma_i$ untuk setiap $i\in I$.
Tunjukkan bahwa
$\|\sup_{i\in I}|X_i|\|_p\le\Bover{p}{p-1}\|X\|_p$.
%275Yd, 275D

\spheader 275Yg Misalkan $(\Omega,\Sigma,\mu)$ adalah ruang probabilitas,
dengan pelengkapan $(\Omega,\hat\Sigma,\hat\mu)$, dan misalkan
$\sequencen{\Sigma_n}$ adalah barisan subaljabar-sigma dari $\hat\Sigma$ yang
tak menurun. Misalkan $\sequencen{X_n}$ adalah barisan fungsi bernilai real
yang terintegralkan terhadap $\mu$ sedemikian sehingga
$\dom X_n\in\Sigma_n$ dan $X_n$ terukur terhadap $\Sigma_n$ untuk setiap
$n\in\Bbb N$. Kita mengatakan bahwa $\sequencen{X_n}$ adalah suatu
{\bf submartingal yang teradaptasi terhadap $\sequencen{\Sigma_n}$} jika
$\int_EX_{n+1}\ge\int_EX_n$ untuk setiap $n\in\Bbb N$ dan setiap
$E\in\Sigma_n$. Buktikan versi 275D, 275F, 275G, 275Xf untuk submartingal.
%275G

\spheader 275Yh Misalkan $\sequencen{X_n}$ adalah suatu martingal, dan
$\phi:\Bbb R\to\Bbb R$ fungsi konveks. Tunjukkan bahwa
$\sequencen{\phi(X_n)}$ adalah suatu submartingal. \Hint{233J.}
Periksa kembali bagian (b-ii) dari bukti 275F berdasarkan fakta ini.
%275F, 275Yg, 275G

\spheader 275Yi Misalkan $\sequencen{X_n}$ adalah barisan variabel acak
nonnegatif yang saling bebas dan semuanya berekspektasi $1$. Tetapkan
$W_n=X_0\times\ldots\times X_n$ untuk setiap $n$. (i) Tunjukkan bahwa
$W=\lim_{n\to\infty}W_n$ terdefinisi hampir di mana-mana. (ii) Tunjukkan
bahwa $\Expn(W)$ bernilai $0$ atau $1$. \Hint{andaikan $\Expn(W)>0$.
Tetapkan $Z_n=\lim_{m\to\infty}X_n\times\ldots\times X_m$. Tunjukkan bahwa
$\lim_{n\to\infty}Z_n=1$ ketika $0<W<\infty$, dan karenanya hampir di
mana-mana berdasarkan hukum nol-satu, sedangkan $\Expn(Z_n)\le 1$ berdasarkan
lemma Fatou, sehingga $\lim_{n\to\infty}\Expn(Z_n)=1$, sedangkan
$\Expn(W)=\Expn(W_n)\Expn(Z_{n+1})$ untuk setiap $n$.} (iii) Tetapkan
$\gamma=\prod_{n=0}^{\infty}\Expn(\sqrt{X_n})$. Tunjukkan bahwa $\gamma>0$
jika dan hanya jika $\Expn(W)=1$.
\Hint{$\Pr(W_n\ge\bover14\gamma^2)\ge\bover14\gamma^2$ untuk setiap $n$,
sehingga jika $\gamma>0$ maka $W$ tidak mungkin nol hampir di mana-mana;
sedangkan $\Expn(\sqrt{W})\le\gamma$.}
%275G

\spheader 275Yj Misalkan $\sequencen{(\Omega_n,\Sigma_n,\mu_n)}$ adalah
barisan ruang probabilitas dengan produk $(\Omega,\Sigma,\mu)$. Anggap bahwa
untuk setiap $n\in\Bbb N$ kita mempunyai ukuran probabilitas $\nu_n$, dengan
domain $\Sigma_n$, yang kontinu mutlak terhadap $\mu_n$, dengan turunan
Radon-Nikod\'ym $f_n$, dan anggap bahwa
$\prod_{n=0}^{\infty}\int\sqrt{f_n}d\mu_n>0$. Misalkan $\nu$ adalah produk
dari $\sequencen{\nu_n}$. Tunjukkan bahwa $\nu$ adalah ukuran integral tak
tentu terhadap $\mu$, dengan turunan Radon-Nikod\'ym $f$, dengan
$f(\pmb{\omega})=\prod_{n=0}^{\infty}f_n(\omega_n)$ untuk $\mu$-hampir setiap
$\pmb{\omega}=\sequencen{\omega_n}$ dalam $\Omega$.
\Hint{gunakan 275Yi untuk menunjukkan bahwa $\int fd\mu=1$.}
%275G, 275Yi

\spheader 275Yk Misalkan $\sequencen{p_n}$ adalah barisan dalam $[0,1]$.
Misalkan $\mu$ adalah ukuran biasa pada $\{0,1\}^{\Bbb N}$ (254J) dan $\nu$
produk $\sequencen{\nu_n}$, dengan $\nu_n$ adalah ukuran probabilitas pada
$\{0,1\}$ yang didefinisikan dengan menetapkan $\nu_n\{1\}=p_n$. Tunjukkan
bahwa $\nu$ adalah ukuran integral tak tentu terhadap $\mu$ jika dan hanya
jika $\sum_{n=0}^{\infty}|p_n-\bover12|^2<\infty$.
%275Yj, 275G

\spheader 275Yl
Temukan martingal $\sequencen{X_n}$ sedemikian sehingga barisan distribusi
$\nu_{X_n}$ (271C) konvergen untuk topologi samar (274Ld), tetapi
$\sequencen{X_n}$ tidak konvergen dalam ukuran.
%275G added 2002

\spheader 275Ym Misalkan $\sequencen{X_n}$ adalah barisan variabel acak
bernilai real yang saling bebas sedemikian sehingga
$\sum_{n=0}^{\infty}X_n$ terdefinisi dalam $\Bbb R$ hampir di mana-mana.
Anggap bahwa terdapat $M\ge 0$ sedemikian sehingga $|X_n|\le M$ hampir di
mana-mana untuk setiap $n$. Tunjukkan bahwa
$\sum_{n=0}^{\infty}\Expn(X_n)$ terdefinisi dalam $\Bbb R$.
\Hint{274Yg, 275G.}
%275G

\spheader 275Yn Misalkan $(\Omega,\Sigma,\mu)$ adalah ruang probabilitas dan
$\sequencen{X_n}$ barisan variabel acak bernilai real yang saling bebas pada
$\Omega$; tetapkan
$E_n=\{\omega:\omega\in\dom X_n$, $|X_n(\omega)|>1\}$,
$Y_n=X_n\times\chi(\Omega\setminus E_n)$ untuk setiap $n$, dan
$Z_n(\omega)=\med(-1,X_n(\omega),1)$ untuk $n\in\Bbb N$ dan
$\omega\in\dom X_n$. Tunjukkan bahwa pernyataan berikut ekuivalen:
(i) $\sum_{n=0}^{\infty}X_n(\omega)$ terdefinisi dalam $\Bbb R$ untuk hampir
setiap $\omega$; (ii) $\sum_{n=0}^{\infty}\hat\mu E_n<\infty$,
$\sum_{n=0}^{\infty}\Expn(Y_n)$ terdefinisi dalam $\Bbb R$, dan
$\sum_{n=0}^{\infty}\Var(Y_n)<\infty$, dengan $\hat\mu$ adalah pelengkapan
$\mu$; (iii) $\sum_{n=0}^{\infty}\hat\mu E_n<\infty$,
$\sum_{n=0}^{\infty}\Expn(Z_n)$ terdefinisi dalam $\Bbb R$, dan
$\sum_{n=0}^{\infty}\Var(Z_n)<\infty$. \Hint{273K, 275Ym.}
(Ini adalah suatu versi dari {\bf Teorema Tiga Deret}.)
%275Ym 275G 273K

\spheader 275Yo Misalkan $(\Omega,\Sigma,\mu)$ adalah ruang probabilitas,
$\sequencen{\Sigma_n}$ barisan subaljabar-sigma dari $\Sigma$ yang tak
menurun dan $\sequencen{X_n}$ barisan variabel acak pada $\Omega$ sedemikian
sehingga $\Expn(\sup_{n\in\Bbb N}|X_n|)$ berhingga dan
$X=\lim_{n\to\infty}X_n$ terdefinisi hampir di mana-mana. Untuk setiap $n$,
misalkan $Y_n$ adalah suatu ekspektasi bersyarat dari $X_n$ pada $\Sigma_n$.
Tunjukkan bahwa $\sequencen{Y_n}$ konvergen hampir di mana-mana menuju suatu
ekspektasi bersyarat dari $X$ pada aljabar-sigma yang dibangkitkan oleh
$\bigcup_{n\in\Bbb N}\Sigma_n$.
%275I %mt27bits

\spheader 275Yp Tunjukkan bahwa 275Yo dapat gagal jika $\sequencen{X_n}$
hanya terintegralkan secara seragam, alih-alih didominasi oleh suatu fungsi
terintegralkan.
%275Yo, 275I  %mt27bits

\spheader 275Yq Misalkan $(\Omega,\Sigma,\mu)$ adalah ruang probabilitas
lengkap, dan $\sequencen{X_n}$ barisan variabel acak yang saling bebas pada
$\Omega$, semuanya dengan distribusi yang sama, dan berekspektasi berhingga.
Untuk setiap $n$, tetapkan $S_n=\bover1{n+1}(X_0+\ldots+X_n)$; misalkan
$\Sigma_n$ adalah aljabar-sigma yang didefinisikan oleh $S_n$ dan
$\Sigma_n^*$ aljabar-sigma yang dibangkitkan oleh
$\bigcup_{m\ge n}\Sigma_m$. Tunjukkan bahwa $S_n$ adalah suatu ekspektasi
bersyarat dari $X_0$ pada $\Sigma_n^*$.
\Hint{anggap setiap $X_i$ terdefinisi di seluruh $\Omega$. Tetapkan
$\phi(\omega)=\sequence{i}{X_i(\omega)}$. Tunjukkan bahwa
$\phi:\Omega\to\BbbR^{\Bbb N}$ adalah \imp\ untuk suatu ukuran produk yang
sesuai pada $\BbbR^{\Bbb N}$, dan bahwa setiap himpunan dalam $\Sigma_n^*$
berbentuk $\phi^{-1}[H]$ dengan $H\subseteq\BbbR^{\Bbb N}$ suatu himpunan
Borel yang invarian terhadap permutasi koordinat dalam himpunan
$\{0,\ldots,n\}$, sehingga $\int_EX_i=\int_EX_j$ kapan pun $i\le j\le n$
dan $E\in\Sigma_n^*$.} Dari sini tunjukkan bahwa $\sequencen{S_n}$ konvergen
hampir di mana-mana. (Bandingkan 273I.)
%275K

\spheader 275Yr Rumuskan dan buktikan versi hasil-hasil dalam bagian ini
untuk martingal yang terdiri atas fungsi bernilai $\Bbb C$ atau $\BbbR^r$,
bukan $\Bbb R$.
}%end of exercises

\cmmnt{
\Notesheader{275} Saya berharap uraian di atas, meskipun sangat disingkat,
telah memberi gambaran tentang kekayaan konsep-konsep yang terlibat, dan akan
menyediakan landasan untuk pembelajaran lebih lanjut. Semua teorema dalam
bagian ini mempunyai implikasi yang luas, tetapi yang benar-benar sangat
diperlukan dalam teori ukuran lanjut ialah 275I, `teorema konvergensi martingal
L\'evy', yang akan saya gunakan dalam pembuktian Teorema Lifting pada Bab 34
dalam jilid berikutnya.

Mengenai waktu henti, saya menyebutkannya sebagian sebagai upaya untuk
menjelaskan lebih jauh kegunaan martingal (lihat 276Ed di bawah), dan sebagian
karena gagasan 275N-275O begitu penting dalam teori probabilitas modern
sehingga, sekadar sebagai pengetahuan umum, Anda perlu menyadari bahwa ada
sesuatu di sana. Saya menambahkan 275P sebagai salah satu hasil baku yang
paling mudah diakses dan dapat diperoleh dengan metode ini.

}%end of notes

\discrpage
